\documentclass[11pt]{article}

\usepackage{fullpage}

\begin{document}

\title{Assembler, LED, and Text-Based Tamagotchi}
\author{Tangina Kader, Amy Showell, Shriya Sistla, Yao Tong}

\maketitle

\section{Assembler Implementation}

Our assembler is composed of three main components: the emulate file, where
instruction fetch and encode is performed; the modules, which consist of 
all instruction-encoding functions; and the data structures. One such data 
structure is the dynamic \texttt{symtable}, a dictionary of labels to addresses. 
The second data structure is the alias table, a dictionary of mnemonics to 
function pointers where the instruction type is handled. 

This assembler utilizes two-pass assembly. The function \texttt{first\_pass} 
loops through the input assembly file, identifying and storing labels in 
a symbol table with their corresponding addresses. Next, the function 
\texttt{second\_pass} loops through the assembly file again, tokenising each 
line to identify the instruction name, replace labels with offsets, and 
extract instruction parameters. The instruction name is then inputted to a 
lookup table to obtain a function pointer to the proper encode function. 
For example, inputting the \texttt{add} instruction to lookup table 
returns a pointer to function \texttt{arith} in module \texttt{data-processing}. 
Each of these functions takes the list of tokens, the number of tokens in 
the list, and an integer pointer where the encoded instruction is to be stored. 
The function returns a status indicating encoding success, either 
\texttt{EXIT\_SUCCESS} or \texttt{EXIT\_FAILURE}. Lastly, the encoded instruction
is converted to binary and written to the specified output file. 
Note that in the second pass, any line consisting of only a label, directive, 
or newline is automatically skipped. 

\section{Blinking LED Implementation}

\section{Text-Based Tamagotchi}

\subsection{Tamagotchi Design}

\subsection{Challenges} 

\subsection{Testing}

\section{Group Reflection}

\section{Individual Reflections}

\subsection{Tangina Kader}

\subsection{Amy Showell}

\subsection{Shriya Sistla}

\subsection{Yao Tong}

\end{document}
